% =====================================================================
%  17.2  Weighted DSU and Rollback DSU              sec:c28-weighted
%  Exports: def:c28-weighted  def:c28-pot  prop:c28-parity
%           ex:c28-pottrace   lst:c28-weighted
%           def:c28-rollback  lst:c28-rollback  prop:c28-rollback
%           fig:c28-rollback  lst:c28-rb-driver
%           prob:c28-impostors  lst:c28-impostors-soln
%  All tables hand-computed on the running tree T (opener, fig:c28-tree).
% =====================================================================
\section{Weighted DSU and Rollback DSU}\label{sec:c28-weighted}

\begin{previouslybox}[Section 17.1 left you with two questions]
\begin{enumerate}[nosep, leftmargin=1.6em]
  \item The plain DSU answers only ``same set or not?''  What if the input
        talks in \emph{relations} -- \(x\) and \(y\) agree, \(y\) and \(z\)
        disagree -- and you must answer ``same side?'' between any two items?
  \item \(\secref{c28-dsu}\)'s warning said path compression is a one-way
        street.  Yet \(\secref{c28-offline}\) will need to try a merge,
        answer some queries, and then \emph{take it back}.  What does a DSU
        look like when undo is legal?
\end{enumerate}
This section answers both -- they are the two \emph{variants} of the same
parent-forest, and both are consumed later: potentials by CF~1594D here,
undo by \secref{c28-offline}'s segment-tree-over-time machine.
\end{previouslybox}

\begin{topicintro}
A \term{weighted DSU} (DSU with \term{potentials}) stores, beside each parent
pointer, one extra number: the \emph{difference} between a node and its
parent.  Climbing to the root accumulates the difference to the
representative, so the structure answers \emph{relative} questions -- same
side, opposite side, exact offset -- not just membership.  The parity
\((0/1)\)-weighted DSU is the algebraic form of bipartiteness.

A \term{rollback DSU} is the opposite trade: it keeps \emph{no} extra
information but gives up path compression, so every union touches exactly two
array cells and can be pushed onto an undo stack.  Snapshots cost \(O(1)\),
restoring undoes exactly the merges taken since.

\textbf{Prerequisites.} \secref{c28-dsu}: \(\coderef{c28-find}/\)-union, the
size rule, the warning that compression is irreversible.  Parity
\((0/1)\) arithmetic mod~2 from Chapter~14's games.

\textbf{After this section you can}
\begin{itemize}[nosep]
  \item carry a potential through \code{find} and \code{unite} without losing
        or double-flipping it; \hfill\lv{4}
  \item reduce ``\(A\) says \(B\) is the same/opposite kind'' constraints to a
        parity DSU and extract the maximum-side count per component; \hfill\lv{4}
  \item write a rollback DSU and explain why compression and undo are
        mutually exclusive; \hfill\lv{4}
  \item feed a rollback DSU through add/undo blocks in \(O(\log n)\) per
        merge. \hfill\lv{5}
\end{itemize}
\end{topicintro}

\begin{recognitionbox}[Pattern card -- ``relations, or merges that must be undone'']
\reckey{``person \(a\) claims \(b\) is the same kind / opposite kind as
\(a\)''; bipartiteness under streamed constraints; difference equations
\(x_u - x_v = w\); ``try this batch of unions, answer, then rewind'';
divide-and-conquer over time with edge lifetimes.}
\recuse{
  \item the constraint links \emph{two items} (same/different/exact offset)
        and queries ask about a pair -- weighted DSU, \defref{c28-weighted},
        CF~1594D here;
  \item unions arrive in \emph{blocks} that are later retracted (binary
        search on time, segment tree over time) -- rollback DSU,
        \defref{c28-rollback}, used by \secref{c28-offline};
  \item both at once: \emph{weighted} rollback is legal -- keep the
        potential, still drop compression (CF~1594D's union pattern survives
        an undo stack untouched).
}
\recnotuse{
  \item only membership matters and nothing is undone -- stay with the plain
        DSU of \secref{c28-dsu}; the potential is dead weight;
  \item the graph's \emph{edges} are the objects and you can simply 2-colour
        per component with BFS -- offline in one pass, plain BFS beats both
        variants for pure bipartiteness (CSES~1668);
  \item deletions must be answered \emph{online} -- no DSU variant survives
        that; it is a link-cut tree, outside this book.
}
One-line test: \emph{is every operation ``relate two items'' (forward) or
``undo the last relates'' (backward)?  Potentials answer the first, a stack
answers the second -- and neither allows path compression to be forgotten
about.}
\end{recognitionbox}

% ---------------------------------------------------------------
\subsection{Weighted DSU: Potentials to the Root}
% ---------------------------------------------------------------

\begin{definition}{Weighted DSU}{c28-weighted}
A weighted DSU adds an array \(\code{pot}[1..n]\) to the structure of
\defref{c28-dsu}.  The invariant: for every \(v\),
\[
  \code{pot}[v] \;=\; \Phi(v) - \Phi(\code{parent}[v]),
\]
where \(\Phi\) is the hidden \term{potential} (here: which side of a
bipartition \(v\) lives on).  \(\code{find}(v)\) returns the root \emph{and}
accumulates \(\Phi(v)-\Phi(\text{root})\) on the way; a union with required
offset \(w\) (meaning \(\Phi(u)-\Phi(v)=w\)) re-links the roots so that the
invariant keeps holding.
\end{definition}

\begin{insightbox}[Everything lives in one equation]
With roots \(r_u,r_v\) already found and offsets \(p_u=\Phi(u)-\Phi(r_u)\),
\(p_v=\Phi(v)-\Phi(r_v)\), imposing \(\Phi(u)-\Phi(v)=w\) while attaching
\(r_v\) under \(r_u\) forces
\[
  \code{pot}[r_v] \;=\; \Phi(r_v)-\Phi(r_u) \;=\; w - p_u + p_v .
\]
Derive it once, and the code in \(\coderef{c28-weighted}\) writes itself.
Attach the \emph{other} way and the sign flips -- this single swap is the
most common wrong answer in the whole family (see the traps below).
\end{insightbox}

\begin{definition}{Parity DSU}{c28-pot}
The \term{parity DSU} is the weighted DSU with potentials in
\(\{0,1\}\) and all arithmetic modulo~2 (XOR).  \(\Phi(v)=0\) or \(1\) names
\(v\)'s side; two items are \emph{same-side} iff their potentials to the
common root are equal, \emph{opposite} iff the potentials differ.
``same kind'' contributes offset \(0\), ``opposite kind'' contributes \(1\).
\end{definition}

\begin{proposition}{Parity DSU \(=\) bipartiteness test}{c28-parity}
Processing constraints ``\(u,v\) same'' \(\big(\Phi(u)\oplus\Phi(v)=0\big)\)
and ``\(u,v\) opposite'' \(\big(\Phi(u)\oplus\Phi(v)=1\big)\): a contradiction
appears exactly when some cycle demands a vertex differ from itself, i.e.\
when the constraint graph is not bipartite; otherwise each component's two
sides \(A,B\) are maintained in \(\code{gp}[r][0],\code{gp}[r][1]\) and every
``same/opposite?'' query is answered in \(\bigO{\alpha(n)}\).
\end{proposition}
\begin{proofbox}
A graph is bipartite iff no odd cycle exists.  Fix a component and any root
\(r\); the DSU has assigned each \(v\) a value \(\Phi(v)\) by XOR-ing the
edge parities along \emph{some} path \(v\leadsto r\).  If the graph is
bipartite, every path between two vertices has the same parity (all cycles
even), so \(\Phi\) is well defined and side-by-\(\Phi\) is a valid
2-colouring; the \code{gp} counts track the two colour classes when a union
merges them (matching sides if the offset is~0, crossed if it is~1).  If an
odd cycle closes, the two paths from a vertex around the cycle disagree, and
the ``same root, wrong parity'' test fires -- exactly the contradiction.
\end{proofbox}

\begin{lstlisting}[caption={find and unite with potentials (parity shown; any abelian group works by replacing xor with +/-)}, label={lst:c28-weighted}]
int find(int v){                       // returns root, fixes pot[v] to root
    if (parent_[v] == v) return v;
    int r = find(parent_[v]);
    pot[v] ^= pot[parent_[v]];         // Phi(v)-Phi(old parent)
    return parent_[v] = r;             //   + (old parent - root) = to root
}
bool unite(int u, int v, int w){       // impose Phi(u) xor Phi(v) = w
    int ru = find(u), rv = find(v);
    if (ru == rv)                      // both potentials are now to ru
        return (pot[u] ^ pot[v]) == w; // consistent? else: contradiction
    int x = w ^ pot[u] ^ pot[v];       // needed pot[rv] if rv hangs on ru
    if (sz_[ru] < sz_[rv]) {           // attaching the OTHER way flips the
        parent_[ru] = rv;              // value to store: x -> x xor w-like
        pot[ru] = x;                   // re-derive, or use the symmetric form
        sz_[rv] += sz_[ru];
        gp[rv][0] += gp[ru][0];        // same-side merge: 0-offset pairs
        gp[rv][1] += gp[ru][1];        // crossed merge: swap the halves
    } else {
        parent_[rv] = ru;
        pot[rv] = x;
        sz_[ru] += sz_[rv];
        gp[ru][0] += gp[rv][0];
        gp[ru][1] += gp[rv][1];
    }
    return true;
}
// same/opposite query:
//   find(u)==find(v) && (pot[u]^pot[v]) == 0/1
\end{lstlisting}

When the offset is always~1 (every constraint says ``opposite''), the two
\code{gp} lines must \emph{cross}: \(\code{gp}[\text{surviving}][0]\mathrel{+}=\)
\(\code{gp}[\text{absorbed}][1]\) and vice versa.  The listing keeps one
\code{x} and merges matching counts because the caller's \(w\) already says
which case it is -- reading \(\code{x}\) as ``the absorbed root's side,
relative to the survivor'' makes both cases identical code.

\begin{example}{Potentials on the running tree, every edge ``opposite''}{c28-pottrace}
Take \(T\) (\(\figref{c28-tree}\)) and declare all nine edges ``opposite
kind'' -- the same reduction as \(\exref{c28-merge}\)'s union order, now with
potentials.  After the nine unions, \(\Phi(v)\) equals \(\depth(v)\bmod 2\),
because each tree edge flips parity exactly once on the path to root~1.
\end{example}
\begin{tracebox}[Parity DSU over T's edges in input order]
\[
\renewcommand{\arraystretch}{1.15}
\begin{array}{c|c||cccccccccc|c}
 & \text{unite} & 1&2&3&4&5&6&7&8&9&10 & \text{component sides}\\ \hline
0 & - & 0&-&-&-&-&-&-&-&-&- & roots: $A=\{1\}$ only\\
1 &(1,2) & 0&1&-&&&&&&& & $A=\{1\}$,\ $B=\{2\}$\\
2 &(1,3) & 0&1&1&&&&&&& & $B=\{2,3\}$\\
3 &(2,4) & 0&1&1&0&&&&&& & $A=\{1,4\}$\\
4 &(2,5) & 0&1&1&0&1&&&&& & $B=\{2,3,5\}$\\
5 &(3,6) & 0&1&1&0&1&1&&&& & $B=\{2,3,5,6\}$\\
6 &(3,7) & 0&1&1&0&1&1&0&&& & $A=\{1,4,7\}$\\
7 &(5,8) & 0&1&1&0&1&1&0&0&& & $A=\{1,4,7,8\}$\\
8 &(7,9) & 0&1&1&0&1&1&0&0&1& & $B=\{2,3,5,6,9\}$\\
9 &(9,10)& 0&1&1&0&1&1&0&0&1&0& \text{final}\\
[2pt]
\end{array}
\]
Row~9 reads $\Phi=(0,1,1,0,1,1,0,0,1,0)$ -- exactly $\depth\bmod 2$:
\[
  A=\{1,4,7,8,10\}\;(\text{even depth}),\qquad
  B=\{2,3,5,6,9\}\;(\text{odd depth}),
\]
$5+5=10$ nodes, both sides of $T$'s bipartition (read the $\depth$ row off
$\figref{c28-tree}$: $0,1,1,2,2,2,2,3,3,4$).  Cross-checks on raw paths, no
compression needed:
$\Phi(5)=\text{pot}[5]\oplus\text{pot}[2]=1\oplus1=0$;
$\Phi(10)=\text{pot}[10]\oplus\text{pot}[9]\oplus\text{pot}[7]\oplus\text{pot}[3]=0\oplus1\oplus0\oplus1=0$.
A \emph{tenth} edge inside $T$, say $(6,8)$, would demand
$\Phi(6)\oplus\Phi(8)=1\oplus0=1$ -- and $6,8$ sit at depths $2,3$, an odd
separation, so the demand is \emph{consistent}: trees are always bipartite,
and the parity DSU never rejects one.  Re-asserting an off-tree edge between
\emph{same}-parity nodes, $(1,5)$ with ``opposite'', is the first real
contradiction: $\Phi(1)\oplus\Phi(5)=0\oplus1=1$ would have to equal~$0$.
\end{tracebox}

% ---------------------------------------------------------------
\subsection{Rollback DSU: Merges You Can Undo}
% ---------------------------------------------------------------

\begin{definition}{Rollback DSU}{c28-rollback}
A \term{rollback DSU} is a plain \(\secref{c28-dsu}\) DSU
\textbf{without path compression}: \(\code{find}\) only climbs, and
\(\code{union}\) attaches by size/rank, writing one record
\((\text{attached root},\ \text{surviving root},\ \text{old sizes})\) onto a
stack.  \(\code{snapshot()}\) returns the current stack height (an
\(\bigO{1}\) integer); \(\code{rollback}(s)\) pops until the height is \(s\),
restoring \(\code{parent}\) and \(\code{sz}\) exactly.
\end{definition}

\begin{warningbox}[Compression and undo cannot coexist]
Path compression rewrites \emph{every} node on the search path -- dozens of
implicit writes with no record.  Undoing a union would then have to undo the
compression, which cannot be attributed to that union at all (later finds
compressed those same nodes under a different root).  Union \emph{by size}
is not an optimisation here, it is the \emph{correctness condition}: it
bounds the untouched tree depth by \(\bigO{\log n}\), so each find stays
\(\bigO{\log n}\) forever and each union touches exactly two cells.
\end{warningbox}

\begin{proposition}{Cost of rollback DSU}{c28-rb-cost}
On \(n\) items, \(\code{find}\) is \(\bigO{\log n}\) worst-case, and
\(\code{unite}\) plus \(\code{rollback}\) are \(\bigO{1}\) per record.  Any
algorithm that performs \(U\) union attempts along a DFS into a recursion
tree of snapshots pays \(\bigO{n + U\log n}\) total -- typically
\(\bigO{\log n}\) DSU-operations per level times \(\bigO{n}\) levels.
\end{proposition}
\begin{proofbox}
Without compression the forest is exactly a union-by-size forest: a node's
depth increases only when its whole tree is hung under a tree at least as
large, so depth doubles each time it happens -- \(\le\log n\) times, so
\(\le\log n\) climb.  Each genuine union writes one parent pointer and one
size; a no-op union writes one sentinel record so a later rollback pops it
symmetrically.  Rollback reverses precisely those fields, so the arrays are
bit-identical to the snapshot state, and the \(\log n\) depth bound is
restored along with every merge.
\end{proofbox}

\begin{lstlisting}[caption={rollback DSU: union by size only, stack of (attached, survivor, survivor's old size)}, label={lst:c28-rollback}]
int stk_[ /* record: attached, survivor, old sz of survivor */ ];
int top = 0;                        // stack height == snapshot token
int find(int v){ while (parent_[v]!=v) v = parent_[v]; return v; }  // NO compression
bool unite(int a, int b){           // every call pushes exactly one record
    a = find(a); b = find(b);
    if (sz_[a] < sz_[b]) { int t=a; a=b; b=t; }     // a = survivor (bigger)
    stk_[top++] = {b, a, sz_[a]};   // record even for no-ops (see below)
    if (a == b) return false;       // already together: no-op marker record
    parent_[b] = a;
    sz_[a] += sz_[b];
    return true;
}
void rollback(int to){              // rewind to snapshot 'to' (a height)
    while (top > to) {
        --top;
        int b = stk_[top][0], a = stk_[top][1], oldsz = stk_[top][2];
        parent_[b] = b;             // b is a root again
        sz_[a] = oldsz;             // survivor's size restored
    }
}
\end{lstlisting}

Two discipline rules make the stack trustworthy:
\textbf{(1) push on every attempt}, including no-ops -- then ``\(k\) unions
after snapshot \(s\)'' always means ``pop \(k\) records'', with no special
case where a union failed; \textbf{(2) snapshot/rollback in strict LIFO
order} -- the usage is a DFS or divide-and-conquer, never a random jump back
three merges and forward again.

\begin{lstlisting}[caption={the snapshot/rollback driver pattern used by secref:c28-offline}, label={lst:c28-rb-driver}]
void solve(int lo, int hi) {          // any recursion over a time/answer axis
    int snap = top;                   // O(1) token: current stack height
    for (each union this block adds) unite(u, v);
    /* ... answer queries that see exactly these unions ... */
    solve(left child);                // deeper blocks add and remove freely
    solve(right child);
    rollback(snap);                   // state here == state at entry
}
\end{lstlisting}

\begin{figure}[H]
\centering
\begin{tikzpicture}[diagram, scale=0.9, transform shape,
    dnode/.style={circle, draw=primaryblue, line width=0.9pt, fill=primaryblue!8,
                  minimum size=6.5mm, font=\small\bfseries, text=primaryblue!60!black},
    rnode/.style={circle, draw=accentorange, line width=1.1pt, fill=accentorange!16,
                  minimum size=6.5mm, font=\small\bfseries, text=accentorange!70!black}]
  % left: two trees, survivor root a (size 3) and root b (size 2), just after unite
  \node[rnode] (a) at (0,0) {$a$};
  \node[dnode] (a1) at (-1.1,-1.2) {};
  \node[dnode] (a2) at ( 1.1,-1.2) {};
  \node[dnode] (b)  at ( 2.6,-0.2) {$b$};
  \node[dnode] (b1) at ( 2.6,-1.6) {};
  \draw[tedge] (a)--(a1) (a)--(a2) (a)--(b) (b)--(b1);
  \node[annot, font=\scriptsize, anchor=north] at (0.7,-2.1)
    {just after $\mathit{unite}$: $b\!\to\!a$, $\mathit{sz}[a]=5$};
  % right: after rollback, b is a root again
  \node[rnode] (ra) at (7,0) {$a$};
  \node[dnode] (ra1) at (5.9,-1.2) {};
  \node[dnode] (ra2) at ( 8.1,-1.2) {};
  \node[rnode] (rb)  at (10.4,0) {$b$};
  \node[dnode] (rb1) at (10.4,-1.4) {};
  \draw[tedge] (ra)--(ra1) (ra)--(ra2) (rb)--(rb1);
  \node[annot, font=\scriptsize, anchor=north] at (7,-2.1)
    {after $\mathit{rollback}(s)$: $\mathit{sz}[a]=3$, $b$ a root};
  \draw[->, accentorange, line width=1.2pt] (4.0,-0.6) --
    node[annot, font=\scriptsize, above] {$\mathit{rollback}(s)$} (5.2,-0.6);
\end{tikzpicture}
\caption{The rollback invariant in miniature (not \(T\) -- a generic step).
A union overwrites exactly two cells, \(\code{parent}[b]\) and
\(\code{sz}[a]\); the record \((b,a,\text{old}\ \code{sz}[a])\) captures both,
so rewinding is a two-field restore, leaving every other pointer and the
\(\bigO{\log n}\) depth bound untouched.  With path compression the same union
would have silently rewritten a whole search path -- unrecoverable from a
single record, which is why compression is forbidden here.  The full
add/undo run on \(T\)'s nine edges is the table below.}
\label{fig:c28-rollback}
\end{figure}

\begin{tracebox}[Add / snapshot / undo on T's nine edges, in input order]
Start ten singletons, stack empty.  Every union resolves both endpoints to
their roots first (no compression), so \(T\)'s growing root~1 (a genuine
singleton is always the smaller side) absorbs each new node; records are
\((\text{attached},\ \text{survivor},\ \text{old }\code{sz})\).
\[
\renewcommand{\arraystretch}{1.15}
\begin{array}{c|l|l|c}
\text{step} & \text{operation} & \text{effect} & \text{stack height} \\ \hline
0 & \mathit{snapshot}\ s_0 & & 0\\
1 & \mathit{unite}(1,2) & (2,1,1)\ \text{pushed},\ \code{sz}[1]=2 & 1\\
2 & \mathit{unite}(1,3) & (3,1,2),\ \code{sz}[1]=3 & 2\\
3 & \mathit{unite}(2,4) & \code{find}(2)=1;(4,1,3),\ \code{sz}[1]=4 & 3\\
4 & \mathit{unite}(2,5) & (5,1,4),\ \code{sz}[1]=5 & 4\\
5 & \mathit{unite}(3,6) & \code{find}(3)=1;(6,1,5),\ \code{sz}[1]=6 & 5\\
6 & \mathit{snapshot}\ s_1 & \text{token}=5 & 5\\
7 & \mathit{unite}(3,7) & \code{find}(3)=1;(7,1,6),\ \code{sz}[1]=7 & 6\\
8 & \mathit{unite}(5,8) & \code{find}(5)=1;(8,1,7),\ \code{sz}[1]=8 & 7\\
9 & \mathit{unite}(7,9) & \code{find}(7)=1;(9,1,8),\ \code{sz}[1]=9 & 8\\
10& \mathit{unite}(9,10)& \code{find}(9)=1;(10,1,9),\ \code{sz}[1]=10 & 9\\
11& \mathit{rollback}(s_1) & \text{pop }4;(10,9,8,7)\to\text{roots},\ \code{sz}[1]=6 & 5\\
12& \mathit{rollback}(s_0) & \text{pop }5\text{ more};\ \text{all singletons} & 0
\end{array}
\]
At step~11 the four merges taken after \(s_1\) unwind in LIFO order
\((9,10)\to(7,9)\to(5,8)\to(3,7)\); nodes \(7,8,9,10\) become roots again and
\(\code{sz}[1]\) returns to~6.  \(\code{find}(4)=1\) but
\(\code{find}(10)=10\): the connectivity that step~10 created has been erased
bit-for-bit -- no cell touched by the first five unions was disturbed.
\end{tracebox}

% ---------------------------------------------------------------
\subsection{Problem Lab: CF 1594D The Number of Imposters}
% ---------------------------------------------------------------

\solvedproblem{c28-impostors}{The Number of Imposters}{CF 1594D}
  {https://codeforces.com/problemset/problem/1594/D}{CF 1700 L4}
  {parity weighted DSU, bipartite components, per-component max side}

\begin{problemstatement}[CF 1594D The Number of Imposters]
In one game there are \(n\) players (besides Theofanis), labelled
\(1..n\); each is either a \textbf{crewmate} (always tells the truth) or an
\textbf{imposter} (always lies).  There are \(m\) comments, each
``\, \(i\) \(j\) \texttt{c}'' meaning: player \(i\) says player \(j\) is a
\texttt{c} (\texttt{crewmate} or \texttt{imposter}).  A set of comments is
\textbf{contradictory} if no assignment of roles makes every comment follow
truthfully/lying as appropriate.  For each test case output the
\textbf{maximum possible number of imposters}, or \(-1\) if the comments
contradict each other.

\textbf{Input.} \(t\) test cases (\(1\le t\le10^4\)); per case: \(n\) and
\(m\) (\(1\le n\le2\cdot10^5\), \(0\le m\le5\cdot10^5\)), then \(m\) lines
\(i\ j\) \texttt{c}.  Sums over all cases: \(\sum n\le2\cdot10^5\),
\(\sum m\le5\cdot10^5\).

\textbf{Output.} One integer per test case.

\textbf{Example (first test case).}
\begin{verbatim}
Input:
3 2
1 2 imposter
2 3 crewmate

Output:
2
\end{verbatim}
Players 2 and 3 can be the imposters here (player~1 a crewmate): the truth-
teller 1 names 2 an imposter, and the liar 2 calls 3 a crewmate -- read the
reduction in \secref{c28-weighted} below.
\end{problemstatement}

\begin{attemptgate}[Attempt it yourself first -- \lv{4}, give yourself 30 minutes]
A fair first attempt: 2-SAT it, or brute-force the \(2^n\) role assignments
for tiny \(n\) and stare at what breaks.  The real question is what a single
comment \emph{means} as a constraint between two players.
\peek{``imposter answers always lie'' sounds self-referential; test the two
possible roles of the \emph{speaker} and see what each forces.}
\nudge{In both cases the speaker and the target land on the \emph{same} or
\emph{opposite} side -- the speaker's own role cancels out.  So every comment
is just ``same kind'' or ``different kind'': the parity DSU of
\defref{c28-pot}.}
\fullhint{Build a parity weighted DSU (\(\coderef{c28-weighted}\)) with
\(0\) for \texttt{crewmate}, \(1\) for \texttt{imposter}; a contradiction is
a same-root wrong-parity union.  Each component ends bipartitioned into sides
\((A,B)\); both colourings are valid, one has \(|A|\) imposters and the other
\(|B|\), and components are independent -- take \(\max(|A|,|B|)\) per
component and sum.  The full build is below.}
\end{attemptgate}

\solstep{Understanding}
Three actors, two roles, and comments \emph{about} roles -- it reads like
self-reference, which invites 2-SAT or case analysis on speakers.  The
constraints link \emph{pairs}, queries concern whole components, and \(n\) is
\(2\cdot10^5\): the shape is a DSU.  What is missing is the DSU's payload:
membership alone cannot tell you \emph{how many} imposters a component can
host -- you need which side of the bipartition each player sits on.  That is
exactly a potential.

\solstep{Key Observations}
\begin{observation}
\textbf{The speaker's role cancels.}
If \(i\) is a crewmate, his comment is true; if an imposter, false -- either
way, ``\(i\) says \(j\) is a crewmate'' holds iff \(i\) and \(j\) have the
\emph{same} role, and ``\(i\) says \(j\) is an imposter'' iff they have
\emph{opposite} roles.  A comment is a pure relation \(\Phi(i)\oplus\Phi(j)\).
\end{observation}
\begin{observation}
\textbf{Consistency is bipartiteness.}
The comments are edges of parity~0/1.  A valid assignment exists iff no
cycle has odd total parity -- i.e.\ iff each component's parity constraints
are consistent.  The parity DSU detects the first violation directly
(\propref{c28-parity}); a violation anywhere means \(-1\) for the case.
\end{observation}
\begin{observation}
\textbf{Each component picks its richer side.}
Once a component is consistent, flipping \emph{all} its roles keeps every
relation satisfied -- so both colourings are legal, and the imposter count
for that component is either \(|A|\) or \(|B|\).  Components never constrain
each other, so the global maximum is \(\sum_{\text{components}}\max(|A|,|B|)\).
Players with no comments are size-1 components: \(\max(1,0)=1\), all
imposters -- which is why \(m=0\) answers \(n\).
\end{observation}

\solstep{The Approach Ladder}
\begin{center}\renewcommand{\arraystretch}{1.25}
\begin{tabular}{@{}c p{0.24\linewidth} p{0.30\linewidth} l c@{}}
\toprule
\textbf{\#} & \textbf{Approach} & \textbf{Idea} & \textbf{Time} & \textbf{Verdict}\\
\midrule
1 & \lv{2}\ enumerate & try all \(2^n\) role strings, check comments &
  \(\bigO{2^n m}\) & \tle \\
2 & \lv{3}\ BFS colour & build parity graph, 2-colour per component, count
  sides & \(\bigO{n+m}\) & \ac \\
3 & \lv{4}\ parity DSU & \(\coderef{c28-weighted}\) online: unite + detect +
  side counts in the roots & \(\bigO{(n+m)\,\alpha(n)}\) & \ac \\
\bottomrule
\end{tabular}\end{center}
Rungs 2 and 3 are the same \(\bigO{n+m}\)-class; the DSU version wins the
chapter slot because it is \emph{online} (comments could stream), because its
per-root \(\code{gp}\) table is the pattern reused in the mixed set, and
because it generalises to non-parity potentials the BFS cannot carry.

\solstep{Think It Through With Me}
\begin{thinkbox}
\deadend{Model each comment as a literal and run 2-SAT}
{Overkill with a fatal flaw: 2-SAT finds \emph{a} satisfying assignment, not
the \emph{maximum-imposter} one; optimising over satisfying assignments is
Max-2-SAT, NP-hard in general.  The structure here is bipartite, not
conjunctive -- exploiting it is the whole problem.}
\deadend{DFS-colour the comment graph, but forget isolated players}
{Sample~5 (\(n=5\), \(m=0\)) then answers \(0\) instead of \(5\): every
unvisited player is its own component with sides \((1,0)\) and contributes
\(\max(1,0)=1\).  A classic wrong-answer on exactly this problem.}
\aha{A comment fixes only a \emph{relation}; relations compose by XOR along
paths; a parity DSU carries the XOR to the root, and the two sides are just
counters attached to the root.}
\end{thinkbox}

\solstep{Solution}
Per test case: reset \(\code{parent},\code{sz},\code{pot},\code{gp}\) for
\(1..n\) (the sums over all cases are bounded, so re-initialising the touched
range is fine).  For each comment \((i,j,\texttt{c})\) call
\(\text{unite}(i,\,j,\,w)\) with \(w=0\) if \texttt{c} is \texttt{crewmate}
else \(w=1\).  If any call returns \texttt{false} (same root, mismatched
parity), the case is contradictory: print \(-1\) and keep reading (do not
break the input).  Otherwise sum \(\max(\code{gp}[r][0],\code{gp}[r][1])\)
over all roots \(r\) (found via \(\code{find}(i)=i\)).

\solstep{Proof of Correctness}
\begin{proofbox}
\textbf{Relations.} By Observation~1, comment \((i,j,\texttt{c})\) is
satisfiable exactly when \(\Phi(i)\oplus\Phi(j)=w_{\texttt{c}}\).
\textbf{DSU invariant.} \(\code{find}\) maintains
\(\code{pot}[v]=\Phi(v)\oplus\Phi(\text{root})\) along the parent chain: when
re-parenting \(v\) under \(r\), \(\code{pot}[v]\) is replaced by the XOR over
the whole path, so the invariant is compression-safe; a union attaches
\(r_v\) with \(\code{pot}[r_v]=w\oplus\code{pot}[u]\oplus\code{pot}[v]\),
which makes \(\Phi(u)\oplus\Phi(v)=w\) computable on any later find --
induction over the operations.
\textbf{Contradiction iff \(-1\).} If a union finds a common root and
\(\code{pot}[u]\oplus\code{pot}[v]\neq w\), the cycle formed by the two paths
to the root plus this edge has odd parity, so no assignment exists; conversely
any unsatisfiable instance has an odd-parity cycle, and when its closing edge
is processed its endpoints already share a root (all earlier cycle edges were
merged) with a differing path parity, which the test detects
(\propref{c28-parity}).
\textbf{Optimality of the count.} On a consistent component, any valid
assignment must place every vertex at potential \(\Phi(v)\oplus b\) for one
global bit \(b\) (paths from the root force it), so exactly two assignments
exist per component, with \(|A|\) resp.\ \(|B|\) imposters; \(\max\) is
optimal, and summing over components is valid because they are independent.
\textbf{Termination/AC.} All operations are \(\bigO{\alpha(n)}\) amortised,
totals within the \(2\cdot10^5 / 5\cdot10^5\) bounds.
\end{proofbox}

\solstep{Implementation}
\begin{lstlisting}[caption={CF 1594D full solution (parity weighted DSU)}, label={lst:c28-impostors-soln}]
#include <bits/stdc++.h>
using namespace std;
const int MX = 200005;
int parent_[MX], sz_[MX], pot[MX];    // pot[v] = Phi(v) xor Phi(parent[v])
int gp[MX][2];                        // per root: |side0|, |side1|
int find(int v){
    if (parent_[v] == v) return v;
    int r = find(parent_[v]);
    pot[v] ^= pot[parent_[v]];        // re-base potential to the root
    return parent_[v] = r;
}
bool unite(int u, int v, int w){      // impose Phi(u) xor Phi(v) = w
    int ru = find(u), rv = find(v);   // pot[u], pot[v] now relative to roots
    if (ru == rv) return ((pot[u] ^ pot[v]) == w);
    int x = w ^ pot[u] ^ pot[v];      // Phi(ru) xor Phi(rv): the root offset
    if (sz_[ru] < sz_[rv]) {          // attach ru under rv at offset x
        parent_[ru] = rv; pot[ru] = x;
        sz_[rv] += sz_[ru];           // ru's side-(x) nodes land on rv side 0,
        int a = gp[rv][0] + gp[ru][x];// ru's side-(1^x) nodes on rv side 1
        int b = gp[rv][1] + gp[ru][1 ^ x];
        gp[rv][0] = a; gp[rv][1] = b;
    } else {                          // attach rv under ru at the SAME x
        parent_[rv] = ru; pot[rv] = x;
        sz_[ru] += sz_[rv];           // (XOR is symmetric: Phi(ru)^Phi(rv)
        int a = gp[ru][0] + gp[rv][x];// equals Phi(rv)^Phi(ru), so one x
        int b = gp[ru][1] + gp[rv][1 ^ x];  // serves both attach directions)
        gp[ru][0] = a; gp[ru][1] = b;
    }
    return true;
}
int main(){
    int T; scanf("%d", &T);
    while (T--){
        int n, m; scanf("%d%d", &n, &m);
        for (int i = 1; i <= n; i++){
            parent_[i] = i; sz_[i] = 1; pot[i] = 0;
            gp[i][0] = 1; gp[i][1] = 0;   // singleton: one node at potential 0
        }
        bool bad = false;
        for (int k = 0; k < m; k++){
            int i, j; char c[12];
            scanf("%d%d%11s", &i, &j, c);
            int w = (c[0] == 'i');        // imposter -> 1, crewmate -> 0
            if (!unite(i, j, w)) bad = true;   // keep reading, mark bad
        }
        if (bad){ printf("-1\n"); continue; }
        int ans = 0;
        for (int i = 1; i <= n; i++)
            if (find(i) == i) ans += max(gp[i][0], gp[i][1]);
        printf("%d\n", ans);
    }
}
\end{lstlisting}

Two things to internalise from \(\coderef{c28-impostors-soln}\).
\textbf{The side counts cross on an odd offset.} Attaching a root at offset
\(x\) means every node on the absorbed side\(\,p\) re-emerges on the
survivor's side \(p\oplus x\); so the merge reads \(\code{gp}[\cdot][x]\)
and \(\code{gp}[\cdot][1\oplus x]\) rather than matching indices.  Get this
backwards and the count is wrong only \emph{when a component has an odd
number of imposter edges} -- a hard-to-see WA.  \textbf{The same \(x\) serves
both attach directions}, because over \(\{0,1\}\) the XOR
\(\Phi(r_u)\oplus\Phi(r_v)\) is symmetric -- subtraction \emph{is}
addition.  Replace XOR by \(+/-\) (general weighted DSU) and the two branches
need \(+x\) and \(-x\): the sign-flip trap the insightbox flags is real, it
merely hides in the parity case.

\solstep{Trace on the Official Sample}
The five-case official sample prints \(\texttt{2 4 -1 2 5}\).
\begin{tracebox}[CF 1594D case 1: n=3, comments (1 2 imposter), (2 3 crewmate)]
Init \(\code{parent}=(1,2,3)\), \(\code{pot}=(0,0,0)\),
\(\code{gp}[i]=(1,0)\) each (node \(i\) sits at potential~0 over itself).
\textbf{unite(1,2,1)}: roots \(1,2\); \(x=1\oplus0\oplus0=1\); size tie keeps
\(1\) the survivor, \(2\!\to\!1\), \(\code{pot}[2]=1\).
Cross-merge by \(x=1\): \(\code{gp}[1][0]=1+\code{gp}[2][1]=1+0=1\),
\(\code{gp}[1][1]=0+\code{gp}[2][0]=0+1=1\) -- \(\code{gp}[1]=(1,1)\),
so side0\(=\{1\}\), side1\(=\{2\}\).
\textbf{unite(2,3,0)}: \(\text{find}(2)=1\) leaves \(\code{pot}[2]=1\);
\(x=0\oplus\code{pot}[2]\oplus\code{pot}[3]=0\oplus1\oplus0=1\);
\(3\!\to\!1\), \(\code{pot}[3]=1\).
Cross-merge: \(\code{gp}[1][0]=1+\code{gp}[3][1]=1+0=1\),
\(\code{gp}[1][1]=1+\code{gp}[3][0]=1+1=2\) -- \(\code{gp}[1]=(1,2)\).
One root \(1\): \(\max(1,2)=\boxed{2}\) (side1 \(=\{2,3\}\) are the
imposters -- the official note's ``players 2 and 3'').
\end{tracebox}
\begin{tracebox}[CF 1594D case 3 (contradiction) and case 5 (empty)]
\textbf{Case 3}, \(n=2\): \((1\ 2\ \text{imposter})\) then
\((2\ 1\ \text{crewmate})\).  First union: \(2\!\to\!1\),
\(\code{pot}[2]=1\).  Second: \(\text{find}(2)=\text{find}(1)=1\), so we test
consistency -- \(w=0\) (crewmate) but \(\code{pot}[2]\oplus\code{pot}[1]=1\oplus0=1\neq0\);
\(\text{unite}\) returns \texttt{false}, \(\mathit{bad}=\)true \(\to\)
\(\boxed{-1}\).  \textbf{Case 5}, \(n=5,m=0\): no unions, five singleton roots,
each contributing \(\max(1,0)=1\) -- \(\boxed{5}\), the isolated-player trap
the thinkbox flagged.  Cases 2 and 4 give \(\boxed{4}\) and \(\boxed{2}\) by
the same per-component \(\max\)-of-sides rule.
\end{tracebox}

\solstep{Complexity}
\begin{complexitybox}
\(\sum n\le2\cdot10^5\), \(\sum m\le5\cdot10^5\) across the file: every
comment is one \(\bigO{\alpha(n)}\) \code{unite}, every final scan one
\code{find} per node -- total \(\bigO{(n+m)\,\alpha(n)}\) per case, linear
overall.  Memory \(\bigO{n}\) (\(\code{parent},\code{sz},\code{pot}\) and the
two-column \(\code{gp}\)).  The reset loop over \(1..n\) is charged to the
bounded \(\sum n\).  \textbf{AC.}
\end{complexitybox}

% ---------------------------------------------------------------
\subsection{Common Pitfalls}
% ---------------------------------------------------------------

\begin{warningbox}[Trap 1: path compression inside a rollback DSU]
The two variants are mutually exclusive.  A rollback DSU \emph{must not}
compress \(\secref{c28-dsu}\)-style; only union-by-size keeps each tree
\(\bigO{\log n}\) deep so the parent pointers you saved on the stack still
describe the merge.  Mixing them silently corrupts the undo history -- see
\(\secref{c28-offline}\) for the correct driver.
\end{warningbox}

\begin{warningbox}[Trap 2: wrong-offset side merge (the WA that only bites sometimes)]
When the offset \(x=1\), the absorbed root's side counts swap before adding
(\(\code{gp}[\mathit{surv}][0]+=\code{gp}[\mathit{abs}][1]\)).  Adding
matching indices always gives a bipartite-looking answer that is right on
even components and wrong on odd ones -- a sample-passing, judge-failing bug.
\end{warningbox}

\begin{warningbox}[Trap 3: no-op union skips the stack push]
In the rollback DSU, push a record even when \(\code{find}(u)=\code{find}(v)\).
Otherwise ``rollback to snapshot \(s\)'' pops a different number of merges
than were taken, and the array state desynchronises from the recursion.
\end{warningbox}

% ---------------------------------------------------------------
\subsection{Practice Problems}
% ---------------------------------------------------------------

\begin{stuckbox}
\begin{itemize}[nosep, leftmargin=1.4em]
  \item Is your potential measured to the \emph{root}, and did you re-base it
        inside \code{find} before using it in a query?
  \item On the swap branch, did the \(\code{gp}\) side counts \emph{cross}
        when the offset is odd?
  \item Rollback: did every union attempt push a record, no-op included?
  \item Did you handle \(m=0\) (all singleton, every player an imposter)?
\end{itemize}
\end{stuckbox}

\begin{practiceset}
  \practice{Desk Disorder}{CF 859E}
    {https://codeforces.com/problemset/problem/859/E}
    {CF 2100 L5}
    {each of \(2n\) name-tags claims one of \(n\) desks, and two tags pinned to
     the same desk force them apart -- a consistency question about \emph{which
     desk} a tag sits at.  Model it as a weighted DSU whose per-set payload is
     the implied desk offset (\(\defref{c28-weighted}\)): union the two tags of a
     worker with a \(\pm1\) potential, and the answer is the product over
     components of (number of valid shifts).  The exact ``carry one number per
     root'' machinery of \(\coderef{c28-weighted}\), with no bipartite colouring
     shortcut available.}
  \practice{Largest Component Size by Common Factor}{LC 952}
    {https://leetcode.com/problems/largest-component-size-by-common-factor/}
    {LC Hard L4}
    {numbers join when they share a factor \(>1\).  Union each number with its
     prime factors (a DSU over \(n + \text{primes}\)), then read the largest
     \(\code{sz}\) at each root -- weighted DSU's ``payload at the root'' idea
     with a plain membership weight.}
\end{practiceset}

\paragraph{Back-cites and forwards.}
\(\secref{c28-dsu}\) -- both variants here are that forest with one field
added (potential) or one habit removed (compression).  Chapter~14
(\chref{c27}) -- parity is the \(\operatorname{mex}\)/XOR thinking of its
Grundy fold.  Forward: the rollback DSU is the engine
\(\secref{c28-offline}\) drives with a segment tree over time.

% ---------------------------------------------------------------
\begin{recallbox}[Section 17.2 in 10 seconds]
\textbf{Weighted DSU} carries \(\code{pot}[v]=\Phi(v)-\Phi(\text{parent})\);
\code{find} re-bases it to the root, and a union at offset \(w\) stores
\(x=w\oplus\code{pot}[u]\oplus\code{pot}[v]\) on the attached root (XOR:
same \(x\) both ways; \(+/-\): signs differ).  Parity DSU \(=\)
bipartiteness; per component the answer is \(\max(|A|,|B|)\).  \textbf{Rollback
DSU} drops compression, unites by size, logs two cells per merge on a stack,
and \code{snapshot}/\code{rollback} are \(\bigO{1}\)/\(\bigO{\text{pops}}\).
Compression and undo never coexist.
\selftest{
  \item State the potential invariant and the re-base line in \code{find}.
  \item Why is the same \(x\) correct on both attach branches -- and what
        breaks when the group is \((\mathbb{Z},+)\) instead of \(\{0,1\}\)?
  \item CF~1594D: why \(m=0\) answers \(n\), not \(0\)?
  \item Give the two rules that make a rollback stack trustworthy.
  \item What is \(\code{find}\)'s cost in a rollback DSU, and why?
}
\end{recallbox}
